\documentclass[11pt]{article}
\usepackage[margin=1in]{geometry}
\usepackage{amsmath,amssymb,amsthm,mathtools,esint}
\usepackage{graphicx}
\usepackage{microtype}
\usepackage[hidelinks]{hyperref}

\newtheorem{theorem}{Theorem}
\newtheorem{lemma}{Lemma}
\newtheorem{corollary}{Corollary}
\newcommand{\E}{\mathbb E}
\newcommand{\BMO}{\mathrm{BMO}}
\newcommand{\cO}{\mathcal O}

\title{A binomial sharpening of the cube--BMO\\rearrangement bound}
\author{Autonomous proof packet for arXiv:2011.09111}
\date{11 August 2026}

\begin{document}
\maketitle

\noindent\textbf{Status.} Partial result, likely valid.  The best constants
and the dimension-free question remain open.

\section{Source question}

Burchard, Dafni, and Gibara prove that every rearrangeable
$f\in\BMO(\mathbb R^n)$ satisfies
\begin{equation}\label{eq:source}
 \|f^*\|_{\BMO(\mathbb R_+)}
 \le 2\bigl(1+2\sqrt{n-1}\bigr)\|f\|_{\BMO(\mathbb R^n)}.
\end{equation}
On page 3 they ask for the best constants $C_n$ and whether the constants can
be dimension-free for $n>1$ \cite{BDG}.

\begin{center}
\includegraphics[width=0.96\textwidth]{figures/open_problem_crop.png}
\end{center}

For a finite-measure set $S$, write
\[
 f_S=\frac1{|S|}\int_S f,
 \qquad
 \cO(f,S)=\frac1{|S|}\int_S|f-f_S|.
\]
The usual BMO seminorm is the supremum of this quantity over axis-parallel
cubes.

\section{Idea of the proof}

The $\sqrt n$ term in \eqref{eq:source} enters at one precise point.  A false
cube with $m$ long directions is tiled by $2^m$ ordinary cubes.  Their means
form a function $g$ on the Bernoulli cube $\{0,1\}^m$, and neighboring values
of $g$ differ by at most $4\|f\|_{\BMO}$.  The source applies a variance bound
and Cauchy--Schwarz to obtain
$\E|g-\E g|\le 2\sqrt m\|f\|_{\BMO}$.

The key observation is that the Doob differences of $g$ are symmetric
two-point variables with predictable amplitudes.  For an even convex loss,
backward conditioning can only increase its expectation when every amplitude
is replaced by its maximum.  Thus the exact comparison variable is a
Rademacher sum, not a Gaussian second-moment upper bound.  Its $L^1$ norm is
the mean absolute deviation of a symmetric binomial variable and is smaller
by the asymptotic factor $\sqrt{2/\pi}$.

A separate geometric lemma says that the oscillation on a subset occupying
fraction $\alpha$ of a set is at most $\alpha^{-1}$ times the oscillation on
the containing set.  Circumscribing a false cube supplies a second bound;
taking the better bound for each number of long directions gives the stated
hybrid constant.

\section{Sharp Bernoulli concentration}

\begin{lemma}[convex domination]\label{lem:convex}
Let $X_1,\dots,X_m$ be independent unbiased Bernoulli variables, and let
$g=g(X_1,\dots,X_m)$.  Suppose changing coordinate $i$ while holding all
others fixed changes $g$ by at most $a_i$.  If $\varepsilon_1,\dots,
\varepsilon_m$ are independent uniform signs, then for every even convex
$\Phi:\mathbb R\to\mathbb R$ for which the expectations exist,
\[
 \E\Phi(g-\E g)
 \le
 \E\Phi\!\left(\frac12\sum_{i=1}^m a_i\varepsilon_i\right).
\]
In particular,
\begin{equation}\label{eq:l1dom}
 \E|g-\E g|
 \le \E\left|\frac12\sum_{i=1}^m a_i\varepsilon_i\right|.
\end{equation}
The bound is sharp when $g=\sum_i a_iX_i$.
\end{lemma}

\begin{proof}
Let $M_i=\E[g\mid X_1,\dots,X_i]$ be the Doob martingale.  Conditional on
$X_1,\dots,X_{i-1}$, the two possible values of $M_i$ have difference
$c_i=c_i(X_1,\dots,X_{i-1})$ with $|c_i|\le a_i$: averaging over the future
coordinates preserves the bounded-difference hypothesis.  Hence, after
identifying $2X_i-1$ with $\varepsilon_i$,
\[
 g-\E g=\frac12\sum_{i=1}^m
 c_i(\varepsilon_1,\dots,\varepsilon_{i-1})\varepsilon_i.
\]

For fixed $s$, convexity implies that
\[
 t\longmapsto \frac{\Phi(s+t)+\Phi(s-t)}2
\]
is nondecreasing for $t\ge0$.  Condition first on
$\varepsilon_1,\dots,\varepsilon_{m-1}$ and replace $|c_m|$ by $a_m$.
After averaging in $\varepsilon_m$, the resulting function of the partial
sum is again even and convex.  Repeat backward for $m-1,m-2,\dots,1$.
This gives the claimed domination.  Taking $\Phi(t)=|t|$ proves
\eqref{eq:l1dom}; constant amplitudes give equality.
\end{proof}

Let $B_m\sim\operatorname{Bin}(m,\tfrac12)$ and put
\begin{equation}\label{eq:beta}
 \beta_m:=\E|B_m-m/2|
 =\frac{m}{2^m}\binom{m-1}{\lfloor(m-1)/2\rfloor}.
\end{equation}
The identity follows by splitting the binomial sum at its mean and using
$k\binom mk=m\binom{m-1}{k-1}$.  It also shows
$\beta_{2r-1}=\beta_{2r}$ and that $(\beta_m)$ is nondecreasing.  Stirling's
formula gives
\begin{equation}\label{eq:betaasymp}
 \beta_m=\sqrt{\frac{m}{2\pi}}+O(m^{-1/2}).
\end{equation}

\section{Application to false cubes}

A false cube in $\mathbb R^n$ has side lengths $2s$ in $m$ coordinate
directions and $s$ in the other $n-m$ directions.  Denote the family of all
false cubes by $\mathcal W$.

\begin{theorem}[binomial false-cube comparison]\label{thm:false}
If $R$ is a false cube with $m<n$ long directions, then
\begin{equation}\label{eq:false}
 \cO(f,R)\le (1+4\beta_m)\|f\|_{\BMO(\mathbb R^n)}.
\end{equation}
Consequently,
\begin{equation}\label{eq:W}
 \|f\|_{\BMO_{\mathcal W}}
 \le (1+4\beta_{n-1})\|f\|_{\BMO}.
\end{equation}
The term $4\beta_m$ in \eqref{eq:false} is optimal as a consequence of only
the individual neighboring-cube bounds employed below.
\end{theorem}

\begin{proof}
Tile $R$ by its $2^m$ congruent cubes $Q(\nu)$, indexed by
$\nu\in\{0,1\}^m$, and set $g(\nu)=f_{Q(\nu)}$.  The partition estimate in
\cite[Lemma 3.6]{BDG} gives
\[
 \cO(f,R)\le \|f\|_{\BMO}+\E|g-\E g|.
\]
Means on adjacent equal cubes differ by at most $4\|f\|_{\BMO}$
\cite[Lemma 3.7]{BDG}.  Lemma~\ref{lem:convex}, with every
$a_i=4\|f\|_{\BMO}$, yields
\[
 \E|g-\E g|
 \le 2\|f\|_{\BMO}\E\left|\sum_{i=1}^m\varepsilon_i\right|
 =4\beta_m\|f\|_{\BMO}.
\]
This proves \eqref{eq:false}; monotonicity of $\beta_m$ gives \eqref{eq:W}.
The additive Hamming-weight function attains Lemma~\ref{lem:convex}, proving
the final assertion about this information set.
\end{proof}

Combining Theorem~\ref{thm:false} with the dimension-free factor 2 for the
decreasing rearrangement on the bisection basis $\mathcal W$
\cite[Lemma 3.5]{BDG} immediately gives
\begin{corollary}\label{cor:simple}
For $n\ge2$ and every rearrangeable $f\in\BMO(\mathbb R^n)$,
\begin{equation}\label{eq:simple}
 \|f^*\|_{\BMO(\mathbb R_+)}
 \le 2(1+4\beta_{n-1})\|f\|_{\BMO(\mathbb R^n)}.
\end{equation}
In particular the leading coefficient in the source bound is improved by
$\sqrt{2/\pi}$:
\[
 C_n\le \sqrt{\frac{32(n-1)}\pi}+2+O(n^{-1/2}).
\]
\end{corollary}

\section{A geometric hybrid refinement}

\begin{lemma}[subset-to-container oscillation]\label{lem:subset}
Let $E\subset Q$ have $0<|E|<|Q|$, and put $\alpha=|E|/|Q|$.  Then
\[
 \cO(f,E)\le \alpha^{-1}\cO(f,Q).
\]
\end{lemma}

\begin{proof}
Write $u=f_E$, $v=f_{Q\setminus E}$, and
$w=f_Q=\alpha u+(1-\alpha)v$.  Jensen's inequality on $Q\setminus E$ gives
\begin{align*}
 \cO(f,Q)
 &\ge \alpha\fint_E|f-w|+(1-\alpha)|v-w|\\
 &=\alpha\left(\fint_E|f-w|+|u-w|\right)\\
 &\ge\alpha\fint_E|f-u|.
\end{align*}
The last step is the triangle inequality.
\end{proof}

A false cube with $m$ long directions lies in a cube of side $2s$ and
occupies the fraction $2^{m-n}$ of that cube.  Lemma~\ref{lem:subset}
therefore supplies the second estimate
\begin{equation}\label{eq:container}
 \cO(f,R)\le 2^{n-m}\|f\|_{\BMO}.
\end{equation}

\begin{theorem}[hybrid rearrangement bound]\label{thm:hybrid}
For $n\ge2$, define
\[
 K_n=\max_{1\le m\le n-1}
 \min\{1+4\beta_m,\,2^{n-m}\}.
\]
Then every rearrangeable $f\in\BMO(\mathbb R^n)$ satisfies
\begin{equation}\label{eq:hybrid}
 \boxed{\ \|f^*\|_{\BMO(\mathbb R_+)}
 \le 2K_n\|f\|_{\BMO(\mathbb R^n)}.\ }
\end{equation}
\end{theorem}

\begin{proof}
For each false cube choose the better of \eqref{eq:false} and
\eqref{eq:container}, then maximize over its number $m$ of long directions.
Ordinary cubes (the case $m=n$) cost only 1 and do not enlarge $K_n$.
Apply the factor-2 bisection-basis rearrangement theorem from \cite{BDG}.
\end{proof}

For illustration, $2K_n$ for dimensions $2$ through $12$ is
\[
4,\ 6,\ 6,\ 8,\ 8,\ 8,\ 9.5,\ 9.5,\ 10.75,\ 10.75,\ 11.84375.
\]
The clean bound \eqref{eq:simple} already improves the asymptotic leading
constant.  The hybrid estimate is never larger, is often strictly smaller,
and also uses only elementary geometry.

\section{Verification and limitations}

The accompanying script verifies formula \eqref{eq:beta}, monotonicity through
$m=200$, the displayed $K_n$ values, and the comparison with the source
constant.  It also solves the finite linear programs defining the maximum
mean absolute deviation of a 1-Lipschitz function on $\{0,1\}^m$ for
$m\le3$; the optima equal $\beta_m$.  These computations are checks only;
the proof above is analytic.

The result does not decide whether $(C_n)$ is dimension-free and does not
determine any sharp $C_n$ for $n>1$.  Moreover, sharpness of the discrete
lemma shows that a dimension-free upgrade cannot use only separate bounds on
the edges of the cube of subcube means.  It would need additional geometric
compatibility among those means or a different decomposition.

A bounded novelty search on 11 August 2026 covered the run indexes, exact
title and inequality queries, author/citation searches, and combinations of
``false cube'', ``binomial'', ``mean absolute deviation'', and
``dimension-free''.  The later VMO work by the same authors and related
semigroup-BMO results do not sharpen this cube-BMO constant.  No occurrence of
the present application was found.  Novelty confidence is moderate because
Lemma~\ref{lem:convex} is elementary and may be known in martingale language.

\paragraph{Human-review recommendation.}
Check the backward induction in Lemma~\ref{lem:convex}, the factor
$4\beta_m$ in Theorem~\ref{thm:false}, and the normalization in
Lemma~\ref{lem:subset}.  Subject to those checks, classify this as a
substantial partial improvement, not a full resolution.

\begin{thebibliography}{9}
\bibitem{BDG}
A. Burchard, G. Dafni, and R. Gibara,
\emph{Mean oscillation bounds on rearrangements},
Trans. Amer. Math. Soc. 375 (2022), 4429--4444;
arXiv:2011.09111.
\end{thebibliography}

\end{document}
